\section{Background and Motivation}
\subsection{YatCC and Compilation Tasks}
\warn{(Introduce YatCC and Complier)}

\warn{TODO: }

{\bf YatCC} is a one-stop serving platform that supports application development, deployment, execution, and management across heterogeneous computing environments through unified runtime abstractions.

{\bf Compiler construction tasks} further represent a highly representative production-scale workload because they require deep integration across code generation, build systems, debugging workflows, dependency management, and iterative execution feedback. (LLVM)

{\bf Intermediate artifacts and verification}.
Each task produces a well-defined intermediate artifact that serves as both the \warn{evaluation} object for the current task and the input dependency for the next:
Task~1 produces a \emph{token stream}; Task~2 produces an \emph{AST/ASG}; Task~3 produces \emph{LLVM IR}; Task~4 produces \emph{optimized IR}; Task~5 produces \emph{RV64 assembly}.
Scoring is performed by dedicated test suites that execute the agent's code against comprehensive test cases.
Each task's score ranges from 0 to 100, \warn{with a passing threshold of 60\%}.

Importantly, these workloads have been validated through years of practical use in compiler construction education at Sun Yat-sen University, ensuring their authenticity as real engineering tasks rather than synthetic benchmarks.


\begin{table*}[!t]
\centering
\caption{Benchmark comparison. \checkmark = supported, --- = not supported. \emph{Serial} = tasks form a dependency chain; \emph{Exec.} = execution-grounded scoring; \emph{Resurr.} = resurrection mechanism; \emph{Process} = process metrics as first-class output; \emph{Artifact} = intermediate artifact continuity tracked.}
\small
\begin{tabular}{lcccccc}
\toprule
\textbf{Benchmark} & \textbf{Serial} & \textbf{Exec.} & \textbf{Resurr.} & \textbf{Process} & \textbf{Artifact} & \textbf{Tasks} \\
\midrule
SWE-bench~{\cite{jimenez2024swebench}} & --- & \checkmark & --- & --- & --- & 2.3K \\
AgentBench~{\cite{liu2024agentbench}} & --- & \checkmark & --- & --- & --- & 8 envs \\
GAIA~{\cite{mialon2024gaia}} & --- & --- & --- & --- & --- & 466 \\
OSWorld~{\cite{xie2024osworld}} & --- & \checkmark & --- & --- & --- & 369 \\
WebArena~{\cite{zhou2024webarena}} & --- & \checkmark & --- & --- & --- & 812 \\
Terminal-Bench~{\cite{merrill2026terminalbench}} & --- & \checkmark & --- & --- & --- & 50 \\
SkillsBench~{\cite{li2026skillsbench}} & --- & \checkmark & paired & --- & --- & 1.5K \\
RE-Bench~{\cite{rebench2024}} & --- & \checkmark & --- & partial & --- & 7 \\
\midrule
\eva{} & \checkmark & \checkmark & \checkmark & \checkmark & \checkmark & 6 \\
\bottomrule
\end{tabular}
\label{tab:benchmark_comparison}
\end{table*}

\subsection{Production-level Serving}
\warn{(State why YatCC and Complier are best as assessment)}

Applications deployed on YatCC interact with workflow runtimes, containerized execution environments, external toolchains, and resource-managed service layers within a realistic operational setting. Such workloads naturally introduce challenges that are often absent from conventional benchmarks, including long-horizon execution, multi-stage tool invocation, asynchronous task coordination, failure recovery, environment dependency management, and runtime observability. As a result, YatCC enables comprehensive assessment not only of functional correctness, but also of runtime robustness, orchestration efficiency, operational stability, and adaptability under practical deployment conditions. These characteristics make YatCC particularly suitable for evaluating production-oriented systems beyond static benchmark environments.

Unlike conventional programming benchmarks that primarily focus on localized code synthesis, compiler-related tasks involve large codebases, multi-stage execution pipelines, strict correctness requirements, and complex coordination across parsing, optimization, code generation, linking, and runtime validation. Such workloads require continuous interaction with external compilation and testing tools, iterative diagnosis of execution failures, and consistent management of evolving project states over long execution horizons. These properties closely resemble real-world software engineering and systems development environments, making compiler construction workloads an effective proxy for evaluating the practical capability, reliability, and runtime behavior of modern intelligent systems in production settings.

\subsection{Motivation and Challenges}
\warn{(How to enable assessment in production)}

Existing agent benchmarks suffer from three fundamental limitations: (1)~\emph{task isolation}---each task is independent, preventing \zzc{assessment}{measurement} of cross-task dependency effects; (2)~\emph{outcome-\zzc{centric}{only} evaluation}---only final scores are reported, with no visibility into process cost, failure patterns, or recovery behavior; (3)~\emph{infrastructure coupling}---each benchmark requires custom evaluation infrastructure, making cross-benchmark comparison and model-backend matrix evaluation impractical.
\warn{(4) unlimited computational resources, time, and tool call budgets}

% === POSITIONING SUMMARY ===
Table~\ref{tab:benchmark_comparison} positions \eva{} against existing benchmarks across five dimensions.
\eva{} is the only \xwz{one}{benchmark} that combines serial dependency, execution-grounded scoring, resurrection, process metrics, and artifact continuity tracking.
While SWE-bench and SkillsBench share execution-grounding and paired evaluation (respectively), neither addresses serial dependency or intermediate artifact flow.
RE-Bench records partial process metrics but lacks resurrection and artifact tracking.
The combination of these five dimensions enables \eva{} to measure \warn{evaluation} axes that no existing benchmark captures simultaneously.

To effectively assess the practical capability of agentic models in production-grade systems, the following key capabilities are necessitated:

{\bf A unified runtime infrastructure for heterogeneous model \warn{evaluation}}.
Production environments typically involve diverse models, APIs, execution frameworks, workflow engines, and deployment backends. A practical assessment platform should therefore provide a unified infrastructure that enables reproducible and scalable evaluation across heterogeneous models and runtime systems. Such an infrastructure should support standardized service interfaces, workflow orchestration, environment provisioning, dependency isolation, and runtime management, allowing different models and execution frameworks to be evaluated under consistent operational conditions.
The infrastructure should also enable closed-loop assessment pipelines for automated testing, execution tracing, error diagnosis, result aggregation, and replayable experimentation, which helps better capture the evolving behavior of agentic systems under complex and changing production workloads.

{\bf Realistic workloads and multi-dimensional assessment modes}.
Conventional benchmarks often focus on isolated tasks with short execution horizons, which fail to capture the characteristics of real production systems. Production-grade assessment instead requires realistic workloads involving long-running workflows, external tool invocation, iterative execution feedback, environment interaction, failure recovery, and multi-stage task coordination. Moreover, assessment should support multiple evaluation modes, including functional correctness, execution efficiency, robustness under runtime perturbations, adaptability to dynamic environments, and operational scalability under concurrent workloads.

{\bf End-to-end observability and measurable runtime metrics}.
Evaluating agentic systems in production settings requires comprehensive runtime visibility beyond task success rates alone. The assessment framework should provide observable metrics spanning workflow execution, tool interaction, runtime behavior, resource consumption, latency characteristics, failure propagation, and execution stability. Such observability is essential for understanding not only whether a task succeeds, but also how the system behaves under realistic operational conditions, enabling deeper analysis of reliability, orchestration capability, and production readiness.

Together, these requirements highlight that \textbf{assessing agentic models in production systems is fundamentally a systems-level problem rather than merely a benchmark-level evaluation task}. Effective assessment therefore necessitates integrated runtime infrastructure, realistic workload modeling, comprehensive observability, and scalable orchestration mechanisms.

